\section{The Critical Editions and the Official Text}

The twentieth century split the Vulgate's two roles---historical document
and church Bible---into separate enterprises, and it is essential to know
which question each modern edition answers.

\textbf{Recovering Jerome's text.} The critical projects aim at the
earliest attainable Hieronymian wording, not the medieval or Clementine
text. The Oxford Vulgate---\work{Nouum Testamentum\ldots\ Latine secundum
editionem sancti Hieronymi}, begun by Bishop John Wordsworth and
H.~J. White---published the Gospels from 1889 and reached completion,
under successor editors, only in 1954; Codex Amiatinus is its leading
witness.\endnote{Edition identity, parts (1889--98; 1913--41; 1905--54),
and witness structure from Houghton, \work{Latin New Testament}, ch.~7.}
For the Old Testament, Pius X in 1907 entrusted the parallel labor to the
Benedictines; the resulting Roman edition (\work{Biblia Sacra iuxta
Latinam Vulgatam versionem}) ran from Genesis (1926) to completion in
1995, carried from 1933 by the Abbey of St.~Jerome in Rome, founded for
the purpose.\endnote{The 1907 commission is described first-hand by its
chairman, Aidan Gasquet, ``Revision of Vulgate,'' \work{Catholic
Encyclopedia} vol.~15 (1912), checked at New Advent (including the
pope's charge ``to determine as accurately as possible the text of
St.~Jerome's Latin translation''); the abbey's foundation by Pius XI
(apostolic constitution \work{Inter praecipuas}, 15 June 1933) is
recited in \work{Scripturarum thesaurus} (next notes); the completion
date 1995 is from Houghton, ch.~7. First-volume date 1926 as in the
standard edition descriptions.} The one-volume Stuttgart
Vulgate (\work{Biblia Sacra iuxta vulgatam versionem}, ed.\ Robert
Weber, 1969; 5th edition by Roger Gryson, 2007) distills both: its Old
Testament rests on the Roman edition, its New Testament on the Oxford
edition, and it prints Jerome's prefaces---making it the standard
scholarly text, deliberately distinct from the Clementine in thousands
of small ways and in its very goal.\endnote{Houghton, \work{Latin New
Testament}, ch.~7. Alongside the Vulgate editions stands the Beuron
\work{Vetus Latina} project (institute founded 1945 by Bonifatius
Fischer at the Archabbey of Beuron), editing the Old Latin evidence the
Vulgate replaced---the indispensable control on every claim about what
Jerome changed.}

\textbf{Revising the church's text.} The official line ran differently.
The Second Vatican Council's constitution on divine revelation,
\work{Dei Verbum} (1965), n.~22, honored ``the Latin translation known
as the vulgate'' among the Church's ancient versions while directing
that new vernacular translations be made ``especially from the original
texts of the sacred books''---Jerome's own principle, now applied past
Jerome's own Bible.\endnote{\work{Dei Verbum} 22, English text at
vatican.va, checked 2026-07-24. The Catechism of the Catholic Church
does not treat the Vulgate as such; its n.~131 quotes this very
conciliar teaching that ``access to Sacred Scripture ought to be open
wide to the Christian faithful'' (CCC 131, vatican.va English text,
checked 2026-07-24; bounded to the Catechism's article on Scripture,
CCC 101--141).} Meanwhile Paul VI, executing the Council's order to
finish the revision of the Psalter, established on 29 November 1965 a
Pontifical Commission for the New Vulgate: a revision of the whole Latin
Bible ``word for word'' where the Vulgate accurately renders the
original texts ``such as they are found in modern critical editions,''
and prudently corrected where it does not, in Christian biblical Latin.
The Psalter appeared in 1969, the New Testament in 1970--1971, the Old
Testament in 1976--1977; and on 25 April 1979 John Paul II, by the
apostolic constitution \work{Scripturarum thesaurus}, declared the
\work{Nova Vulgata} the ``typical'' edition and promulgated it for use
``especially in the sacred Liturgy.'' A second typical edition with
retouches followed in 1986.\endnote{\work{Scripturarum thesaurus}
(25 April 1979), English text at vatican.va, checked 2026-07-24 (which
also notes the revision was ``quite demanding in certain books of the
Old Testament which Saint Jerome did not touch''); the working norms,
the commission history under Cardinal Bea and Bishop Schick, and the
volume dates are from the \work{Praefatio ad Lectorem} of the
\textit{editio typica altera}, Latin text at vatican.va, checked
2026-07-24. The Nova Vulgata's New Testament took the Stuttgart critical
text as its base (Houghton, ch.~7).}

The \work{Nova Vulgata} is thus a new thing under an old name: an
official Latin version aligned with the modern critical Hebrew and Greek
texts, replacing the Gallican Psalter, and serving as the Latin liturgy's
reference text---a role the Congregation for Divine Worship's
instruction \work{Liturgiam authenticam} (2001) made explicit by
directing that liturgical translation follow the manuscript tradition
the \work{Nova Vulgata} follows for the canonical text, and consult it
``as an auxiliary tool'' to maintain the Latin tradition of
interpretation, while translating from the original
languages.\endnote{\work{Liturgiam authenticam} (28 March 2001),
nn.~24, 37, 41, and 43, English text at vatican.va, checked 2026-07-24.
The instruction's own footnotes cite \work{Scripturarum thesaurus} and
the 1986 \textit{editio typica altera}. Its directives are liturgical
law of their date, cited here as an act of reception, not as a
statement about textual history.} Three ``Vulgates'' therefore now
coexist, and none may be silently substituted for another: the
Clementine (the historical church text, 1592--1979), the Stuttgart
(the scholarly reconstruction of Jerome), and the \work{Nova Vulgata}
(the current typical edition). Which one a claim concerns is the first
question to ask of any modern sentence containing the word
``Vulgate.''\endnote{Class~S synthesis. Even mechanical identification
differs: at Genesis 3:20 the three print \textit{Heva}, \textit{Hava},
and \textit{Eva} respectively (Houghton, ch.~7).}
